\documentclass[10pt,a4paper]{article}
\usepackage[T1]{fontenc}
\usepackage[utf8]{inputenc}
\usepackage{lmodern}
\usepackage[margin=2.0cm]{geometry}
\usepackage{microtype}
\usepackage{booktabs}
\usepackage{tabularx}
\usepackage{array}
\usepackage{xcolor}
\usepackage[hyphens]{url}
\usepackage[colorlinks=true,linkcolor=black,urlcolor=blue!50!black]{hyperref}
\usepackage{enumitem}
\usepackage{fancyvrb}
\setlist{nosep,leftmargin=1.3em}
\setlength{\parindent}{0pt}
\setlength{\parskip}{5pt}
\renewcommand{\arraystretch}{1.25}
\newcolumntype{L}{>{\raggedright\arraybackslash}X}
\newcommand{\code}[1]{\path{#1}}
\fvset{fontsize=\small,frame=single,framesep=6pt,rulecolor=\color{black!35}}

\begin{document}

{\LARGE\bfseries Guide for judges: ``Is the Map Right?''\par}
\vspace{6pt}
{\large Repository: \url{https://github.com/SauravBellam/is-the-map-right}\par}
\vspace{8pt}

This page tells you where every part of the paper comes from and how to check it. Everything named here is in the repository above; all paths are relative to its \code{planA_results/} folder.

\section*{1. Read}

\begin{tabularx}{\linewidth}{@{}lL@{}}
\toprule
What & File\\
\midrule
Paper (PDF) & \code{papers/latex/paper1_is_the_map_right.pdf}\\
Paper (LaTeX source) & \code{papers/latex/paper1_is_the_map_right.tex}\\
Figures (vector / PNG) & \code{papers/figures/*.pdf}, \code{papers/figures/png/*.png}\\
Numbers behind each figure & \code{papers/figures/figure_data.json}\\
\bottomrule
\end{tabularx}

\section*{2. Check the numbers in one command}

\begin{Verbatim}
cd planA_results
python3 papers/verify_paper1_numbers.py
\end{Verbatim}

Standard library only, a few seconds. It re-reads the result files and compares 66 quantitative claims in the paper with the values in the files, printing PASS or FAIL with the source file of each. The saved output is \code{papers/verify_paper1_numbers.out.txt}: \textbf{66 PASS, 0 FAIL, 9 INFO} (INFO lines list values that appear only in figures). The script exits non-zero if any claim fails.

An earlier draft of the paper had three small wording errors that this script caught (\S5.1, a single non-async test class in rich; \S5.4, the httpx graph at 30\%; \S5.6, a rank-correlation bound of 0.17 rather than 0.16). The paper text was corrected; no result changed.

\section*{3. Where each section's results live}

{\small
\begin{tabularx}{\linewidth}{@{}>{\raggedright\arraybackslash}p{0.2\linewidth}LL@{}}
\toprule
Paper section & Result files & Produced by\\
\midrule
\S5.1 Omission (Fig.~1) & \code{vm-a2/results/a1_v2/a1_summary.json}, \code{a1_tasks.csv}, \code{a1_gold_symbols.jsonl}, \code{a1_graphs.csv} & \code{scripts/planA/a1_census.py}\\
\S5.1 Id reconciliation & \code{vm-a2/results/a7/a7_summary.json}, \code{a7_class_dropped.csv}, \code{a7_provided_nodes.csv}, \code{a7_defs_absent.csv} & \code{scripts/planA/a7_id_audit.py}\\
\S5.2 Resolver (Fig.~2) & \code{vm-a2/results/a1_v2/a1_summary.json} (\code{misresolution}) & \code{scripts/planA/a1_census.py}, resolver copy in \code{scripts/planA/common.py}\\
\S5.3 Recall and precision (Table~1, Fig.~3) & \code{vm-a2/results/a2_v3/a2_summary.json}, \code{a2_tasks.csv}, \code{a2_run_status.csv}, \code{a2_dyn_edges.jsonl}, \code{a2_static_unconfirmed.jsonl} & \code{scripts/planA/a2/run_task.py}, \code{trace_plugin.py}, \code{analyze.py}\\
\S5.4 Edge semantics (Figs.~4 and 5) & \code{vm-a2/results/a8/a8_summary.json}, \code{a8_edges.csv}, \code{a8_sample.csv} (80 edges with caller snippets) & \code{scripts/planA/a8_edge_audit.py}; structure classes in \code{papers/make_figures.py}\\
\S5.5 Reachability (Fig.~6) & \code{vm-a2/results/a2_v3/a2_summary.json} (\code{gold_min_hops_from_entry_hist}), \code{a2_gold.jsonl} & \code{scripts/planA/a2/analyze.py} (\code{entry_nodes})\\
\S5.6 Encoder & \code{vm-a2/results/a3/}, \code{vm-a2/results/a3_gemma2/} & \code{scripts/planA/a3_encoder.py}\\
\S5.6 Similarity (Fig.~7) & \code{vm-main/results/a4/a4_summary.json}, \code{a4_relations.csv}, \code{a4_knn.csv} & \code{scripts/planA/a4_probe.py}\\
\S5.6 Retrieval (Fig.~8) & \code{vm-a2/results/a5_v3/a5_summary.json}, \code{a5_tasks.csv} & \code{scripts/planA/a5_retrieval.py}\\
\S6 Released graph and profile & \code{vm-a2/results/a6_v3/graphs_v2/*.json.gz} (one per commit, gzipped), \code{a6_v3/task_profile.csv} & \code{scripts/planA/a6_release.py}\\
Per-task tests and traces & \code{vm-a2/results/a2_runs/<task>/} (status, pytest logs, junit XML, traces, dependency logs) & \code{scripts/planA/a2/run_task.py}\\
\bottomrule
\end{tabularx}
}

\section*{4. Reproduce from scratch}

The competition archive (tasks, graphs, embeddings, snapshots) is not redistributed here; download it from the competition. Then, on a Linux machine with Python 3.13 and network access:

\begin{Verbatim}[fontsize=\footnotesize]
cd planA_results/scripts
python3 planA/run_all.py --data <archive> --out <results> --stages a1,a4,a5 --a1-shards 8 --jobs 8
python3 planA/a2/run_task.py --data <archive> --out <results>/a2_runs --backend venv --jobs 16 --keep
bash ops/rerun_fixes.sh <archive> <results> <cache> 8
bash ops/run_gaps.sh <archive> <results> <cache>
\end{Verbatim}

Measured run times (from the step logs): census 2.4 minutes on 8 shards; the dynamic oracle 9,756 task-seconds in total (median 57\,s per task, so roughly 10 minutes of wall time at 16 parallel jobs, plus up to 25 minutes for a task whose tests hang); every analysis step under 5 minutes. The encoder step needs a Hugging Face token with access to EmbeddingGemma. Host specs and package versions are in \code{provenance/}; script checksums for each pass are in \code{vm-a2/results/scripts_sha256_v3.txt}.

Tests of the pipeline itself (synthetic fixture, no competition data needed):

\begin{Verbatim}
cd planA_results/scripts && python3 -m pytest tests/test_pipeline.py -q
\end{Verbatim}

15 tests pass.

\section*{5. History and caveats worth knowing}

\begin{itemize}
\item The first run (Oct 10) used a wrong id scheme for nested definitions; it was fixed and every affected step rerun. The repository contains only the final outputs (\code{a1_v2}, \code{a2_v3}, \code{a5_v3}, \code{a6_v3}, \code{a7}, \code{a8}), and the paper uses only these.
\item One task (rich\_4006) hung on four test cases; it was recovered by deselecting them and recording them as failed (\code{vm-a2/results/a2_runs/rich_4006/status.json}, \code{hung_tests}).
\item 15 of 129 tasks did not run cleanly in our environment (9 collection errors, 2 import errors, 4 requests tasks with no collectable tests). These reflect our dependency setup, not the tasks.
\end{itemize}

\section*{6. Large files}

The 127 released typed graphs are stored gzipped (\code{vm-a2/results/a6_v3/graphs_v2/*.json.gz}, 662\,MB uncompressed); \code{papers/make_figures.py} reads them as they are. The per-definition census rows are in \code{vm-a2/results/a1_v2/a1_async_audit.csv.gz}.

\end{document}
